\documentclass[10pt,twocolumn,a4paper]{article}
\usepackage[utf8]{inputenc}
\usepackage{mathptmx}
\usepackage{amsmath,amssymb,amsfonts,amsthm}
\usepackage{geometry}
% Ajuste estricto de márgenes para maximizar espacio (1 a 2 páginas)
\geometry{top=1.85cm, bottom=2.15cm, left=1.5cm, right=1.5cm, columnsep=0.65cm}
\usepackage{booktabs}
\usepackage{xcolor}
\usepackage{titlesec}
\usepackage{caption}
\usepackage{enumitem}
\usepackage{microtype}
\usepackage{hyperref}
\hypersetup{
    colorlinks=true,
    linkcolor=blue!75!black,
    citecolor=blue!75!black,
    urlcolor=blue!75!black,
    pdfencoding=auto,
    psdextra
}

\newtheorem{definition}{Definición}
\newtheorem{lemma}{Lema}
\newtheorem{theorem}{Teorema}

% Formato de encabezados estilo IEEE
\titleformat{\section}{\normalsize\bfseries\scshape\filcenter}{\thesection.}{0.5em}{}
\titleformat{\subsection}{\normalsize\bfseries\itshape}{\thesubsection.}{0.5em}{}
\titleformat{\subsubsection}{\small\bfseries}{\thesubsubsection.}{0.5em}{}

\titlespacing*{\section}{0pt}{8pt plus 2pt minus 2pt}{4pt plus 1pt minus 1pt}
\titlespacing*{\subsection}{0pt}{6pt plus 2pt minus 1pt}{3pt plus 1pt minus 1pt}
\titlespacing*{\subsubsection}{0pt}{4pt plus 1pt minus 1pt}{2pt plus 1pt minus 1pt}

\captionsetup[table]{name=TABLA, labelfont={bf,scriptsize}, textfont=scriptsize, labelsep=colon, skip=2pt}
\captionsetup[figure]{name=Fig., labelfont={bf,scriptsize}, textfont=scriptsize, labelsep=period, skip=2pt}

\begin{document}

\title{\vspace{-0.9cm}\textbf{\Large Diseño, Correctitud y Análisis de Complejidad de un Range Tree 2D con Fractional Cascading}}

\author{
    \textbf{Autores:} Thiago Rivera, Wilder Bañon, Tamy Flores, Jesús Niño \\
    \textbf{Proyecto 2: Efecto Cascada} \\
    \textit{CS3014 - Estructuras de Datos Avanzadas} \\
    \textit{Universidad de Ingeniería y Tecnología (UTEC)} - Lima, Perú \\
}
\date{}
\maketitle

\noindent\small\textbf{\textit{Resumen}---En este informe se presenta el diseño, análisis matemático de correctitud y evaluación de complejidad de un \textit{Layered Range Tree} bidimensional optimizado con \textit{Fractional Cascading}. Esta estructura resuelve consultas de conteo ortogonal sobre un conjunto de hasta $3 \cdot 10^5$ puntos en tiempo $O(\log n)$, superando el límite teórico de $O(\log^2 n)$ de los árboles de rangos clásicos. Se formaliza el invariante de los arreglos en cascada, la lógica de preprocesamiento ascendente tipo \textit{MergeSort} y el descenso directo en $O(1)$. Asimismo, se demuestra mediante inducción y recurrencias matemáticas que el uso de memoria y tiempo de construcción se mantienen estrictamente en $O(n \log n)$.}

\vspace{0.4em}
\noindent\small\textbf{\textit{Palabras clave}---Range Tree, Fractional Cascading, Geometría Computacional, Búsqueda Ortogonal.}

\vspace{0.5em}

\section{Introducción}
El problema de búsqueda en rangos ortogonales exige identificar elementos multidimensionales contenidos en hiper-rectángulos de consulta~\cite{deberg}. En su versión bidimensional de conteo, dado un conjunto de puntos $P$ y una caja $R = [x_1, x_2] \times [y_1, y_2]$, se debe retornar $|P \cap R|$~\cite{doc}. 

La solución clásica emplea un \textit{Range Tree 2D}, el cual requiere $O(\log^2 n)$ por consulta debido a búsquedas binarias redundantes en árboles secundarios~\cite{deberg}. Dadas las restricciones de $n, q \le 3 \cdot 10^5$ y un límite de 5 segundos~\cite{doc}, este proyecto implementa un \textit{Layered Range Tree} aplicando \textit{Fractional Cascading} (o \textit{cross-linking}). Esta técnica permite unificar las búsquedas secundarias precalculando punteros fraccionarios, reduciendo el tiempo de consulta a $O(\log n)$~\cite{deberg}.

\section{Arquitectura y Diseño de la Estructura}

\subsection{Subconjuntos Canónicos y Arreglos Secundarios}
Sea $T$ el árbol principal construido sobre la dimensión $X$. Cada nodo $u \in T$ define un subconjunto canónico $P(u)$ que contiene las hojas de su subárbol~\cite{deberg}. En lugar de usar un BST secundario, en un \textit{Layered Range Tree}, $P(u)$ se almacena en un arreglo contiguo \texttt{currentY}, el cual está estrictamente ordenado respecto a la coordenada $Y$~\cite{doc, deberg}. Para minimizar el uso de memoria, solo la raíz de $T$ preserva explícitamente las coordenadas $Y$ de los puntos en \texttt{rootY}.

\subsection{Puentes Estructurales (Fractional Cascading)}
La clave del diseño radica en la conexión de las estructuras secundarias. Cada nodo $u$ mantiene dos arreglos de enteros, \texttt{toLeft} y \texttt{toRight}, mapeados uno a uno con su arreglo \texttt{currentY}~\cite{doc}.

\begin{definition}[Puentes de Cascada~\cite{deberg}]
Para un elemento con índice $i$ en el arreglo \texttt{currentY} del nodo $u$, el valor $\texttt{toLeft}[u][i]$ almacena el menor índice $j$ en el hijo izquierdo tal que $\text{leftY}[j] \ge \texttt{currentY}[i]$. Análogamente se define $\texttt{toRight}[u][i]$ para el hijo derecho.
\end{definition}

Esta arquitectura desacopla el costo de la búsqueda; en lugar de hacer repetidos \textit{binary searches}, un índice en el nodo $u$ se traduce a su índice correspondiente en los nodos hijos en tiempo $O(1)$.

\section{Mecánica de Operaciones}

\subsection{Preprocesamiento Ascendente (\texttt{build})}
La construcción de los arreglos se ejecuta \textit{bottom-up}. En las hojas, \texttt{currentY} contiene un único elemento. En los nodos internos, se combinan los arreglos \texttt{leftY} y \texttt{rightY} de sus hijos mediante una técnica \textit{Two-Pointers} idéntica a la fase de mezcla de \textit{MergeSort}~\cite{doc}. 

Inmediatamente después de la fusión, se computan los puentes \texttt{toLeft} y \texttt{toRight}. Se itera sobre el arreglo recién formado \texttt{currentY} con un puntero $i$, y se avanzan los punteros locales $j$ y $k$ en \texttt{leftY} y \texttt{rightY} respectivamente para registrar las transiciones.

\subsection{Consulta de Conteo (\texttt{query})}
La consulta evalúa el rectángulo $R = [x_1, x_2] \times [y_1, y_2]$ de la siguiente forma~\cite{doc}:
\begin{enumerate}[leftmargin=*]
    \item Se localizan los límites $[ql, qr]$ en el eje $X$ sobre el arreglo de puntos base.
    \item Se ejecuta la \textbf{única} búsqueda binaria sobre el eje $Y$ en el arreglo \texttt{rootY} de la raíz, obteniendo los índices $[y_{low}, y_{high})$. Matemáticamente, esto se traduce en computar el \textit{lower bound} (límite inferior estricto) para $y_1$ y el \textit{upper bound} (límite superior estricto) para $y_2$, delimitando con precisión el intervalo de interés.
    \item Se inicia un recorrido descendente desde la raíz. Si el rango horizontal del nodo $u$ está completamente contenido en $[ql, qr]$, se suma $(y_{high} - y_{low})$ al total y no se desciende más.
    \item Si se debe descender, los nuevos límites para los hijos se obtienen mediante $y_{low}' = \texttt{toLeft}[u][y_{low}]$ y $y_{high}' = \texttt{toLeft}[u][y_{high}]$, realizando la transición al hijo izquierdo en tiempo $O(1)$. Lo mismo aplica al hijo derecho con \texttt{toRight}.
\end{enumerate}

\section{Análisis de Correctitud}

\begin{theorem}[Correctitud de Transición Fraccionaria]
Si el intervalo de índices $[y_{low}, y_{high})$ en el nodo $u$ contiene exactamente los puntos de $P(u)$ cuyas coordenadas Y están en $[y_1, y_2]$, entonces el intervalo $[\texttt{toLeft}[u][y_{low}], \texttt{toLeft}[u][y_{high}])$ delimita exactamente los puntos correspondientes en el hijo izquierdo.
\end{theorem}

\begin{proof}
Procedemos por inducción sobre la profundidad del descenso en el árbol.
\textbf{Caso base:} En la raíz, la búsqueda binaria inicial garantiza que el intervalo $[y_{low}, y_{high})$ acota de manera exacta a los elementos cuyas coordenadas $Y \in [y_1, y_2]$.
\textbf{Hipótesis inductiva:} Supongamos que en un nodo $u$, el intervalo de índices $[y_{low}, y_{high})$ delimita exactamente al subconjunto $P(u) \cap R_Y$, donde $R_Y = [y_1, y_2]$.
\textbf{Paso inductivo:} Sea $v$ el hijo izquierdo de $u$. Al aplicar $y_{low}' = \texttt{toLeft}[u][y_{low}]$, por la Definición 1, $y_{low}'$ es el menor índice en $v$ tal que $\texttt{leftY}[y_{low}'] \ge \texttt{currentY}[y_{low}] \ge y_1$. Simétricamente, $y_{high}' = \texttt{toLeft}[u][y_{high}]$ mapea al primer elemento que excede estrictamente $y_2$. Dado que $P(v) \subseteq P(u)$ y ambos arreglos mantienen un ordenamiento estricto, la correspondencia biunívoca preserva las cotas. Por lo tanto, el nuevo intervalo $[y_{low}', y_{high}')$ delimita exactamente a $P(v) \cap R_Y$.
\end{proof}

\begin{lemma}[Exactitud del Conteo Canónico]
Si un nodo $u$ es canónico respecto al rango de la consulta en $X$ (es decir, $X_u \subseteq [x_1, x_2]$), entonces $y_{high} - y_{low}$ es exactamente el número de puntos en $P(u) \cap R$.
\end{lemma}

\begin{proof}
Al ser $u$ canónico en $X$, todos los puntos en $P(u)$ satisfacen la restricción horizontal. Como el intervalo semiabierto $[y_{low}, y_{high})$ abarca todos los puntos en $P(u)$ que caen en $[y_1, y_2]$ (demostrado por el Teorema 1), y la estructura utiliza indexación desde cero, la diferencia aritmética simple $y_{high} - y_{low}$ determina con exactitud la cardinalidad del conjunto sin necesidad de iteración adicional ni correcciones de desplazamiento (offset).
\end{proof}

\section{Análisis Formal de Complejidad}

\subsection{Complejidad Espacial}
En un árbol perfectamente balanceado, cada punto original aparece exactamente en un nodo por cada nivel de profundidad. Un nodo a profundidad $d$ mantiene su arreglo canónico, cuyos arreglos de puentes \texttt{toLeft} y \texttt{toRight} son de tamaño $|P(u)| + 1$. 
Sea $S(n)$ el espacio total, la suma de los tamaños de los arreglos en cualquier nivel fijo es $O(n)$~\cite{deberg}. Como la altura del árbol base es $\log_2 n$, tenemos:
\begin{equation}\label{eq:espacio}
S(n) = \sum_{d=0}^{\log_2 n} O(n) = O(n \log n)
\end{equation}

\subsection{Tiempo de Preprocesamiento}
La construcción recursiva genera la recurrencia $T(n) = 2T(n/2) + f(n)$. El término $f(n)$ corresponde a la mezcla (Merge) y el cálculo de puentes. Como ambos procedimientos utilizan punteros que iteran secuencialmente sin retroceso, procesar un nodo toma tiempo estrictamente proporcional a su tamaño: $f(n) = O(n)$~\cite{doc}. Evaluando la recurrencia mediante el Teorema Maestro:
\begin{equation}\label{eq:tiempo_build}
T(n) = 2T(n/2) + O(n) \implies T(n) = O(n \log n)
\end{equation}

\subsection{Tiempo de Consulta en el Peor Caso}
El tiempo de consulta $Q(n)$ depende de la búsqueda inicial y el recorrido del árbol. 
\begin{itemize}[leftmargin=*]
    \item \textbf{Fase 1:} $O(\log n)$ para hallar los límites iniciales en $X$ y en $Y$.
    \item \textbf{Fase 2:} En el peor caso, el rango $X$ se descompone en a lo sumo $2 \log_2 n$ subárboles canónicos (nodos delimitadores).
    \item \textbf{Fase 3:} Por Fractional Cascading, evaluar la condición y bifurcar los límites toma tiempo constante $O(1)$ por nodo visitado.
\end{itemize}
Por ende, el tiempo de consulta total es:
\begin{equation}
Q(n) = \underbrace{O(\log n)}_{\text{Búsquedas base}} + \underbrace{O(\log n) \times O(1)}_{\text{Descenso con punteros}} = O(\log n)
\end{equation}
Al ejecutar $q$ consultas, el tiempo global del algoritmo resulta $O(n \log n + q \log n)$, superando ampliamente el margen requerido por el problema ($O(q \log^2 n)$).

\section{Conclusiones}
\begin{enumerate}[leftmargin=*]
    \item Se diseñó e implementó un \textit{Layered Range Tree 2D}, integrando \textit{Fractional Cascading} mediante puentes estáticos calculados \textit{bottom-up}.
    \item El análisis demostró que el cálculo exhaustivo de índices fraccionarios no altera la cota teórica de preprocesamiento, la cual se mantiene en $O(n \log n)$.
    \item La sustitución de operaciones repetidas de búsqueda binaria por simples traslaciones de índices acotó estrictamente la consulta de recuento ortogonal a $O(\log n)$.
    \item La eliminación de estructuras asociadas complejas a favor de arreglos paralelos optimizó la localidad en caché, garantizando un desempeño robusto dentro de los límites de tiempo del sistema de evaluación.
\end{enumerate}

\begin{thebibliography}{99}
\bibitem{doc}
UTEC - Ciencia de la Computación, \textit{Documentos del curso: Proyecto 2 - Efecto Cascada (Instrucciones y código base)}, 2026.
\bibitem{deberg}
M. de Berg, O. Cheong, M. van Kreveld, y M. Overmars, \textit{Computational Geometry: Algorithms and Applications}, 3ra ed. Capítulo 5: Orthogonal Range Searching.
\end{thebibliography}

\end{document}
